% ─────────────────────────────────────────────────────────────
\phantomsection
\chapter{Related Work}
\label{sec:related}

Wearable fall detection has been approached from three main directions: rule-based
threshold detectors, classical machine-learning pipelines built on handcrafted
features, and end-to-end deep-learning models.  Each direction has produced systems
that achieve strong scores on controlled benchmark datasets, yet each carries
structural limitations that become apparent when the goal shifts from benchmark
performance to robust, personalised deployment.  A fourth line of research applies
persistent homology to time-series classification in domains outside of fall detection.
This section surveys each strand, identifies the specific limitations that motivate
the present work, and closes with an explicit positioning statement.

% ─────────────────────────────────────────────────────────────
\section{Threshold-Based Fall Detectors}
\label{subsec:rw-threshold}

\subsection{Method family}

The oldest and most widely deployed family of wearable fall detectors operates
directly on the signal magnitude vector
$\mathrm{SMV}(t)=\|\mathbf{a}(t)\|_2$
and raises an alarm whenever $\mathrm{SMV}(t)$ exceeds a fixed scalar threshold~$T$:
\[
  \widehat y(t) = \mathbf{1}\{\mathrm{SMV}(t) > T\}.
\]
Kangas et al.~\parencite{kangas2008fall} established the archetypal version of this
approach on a waist-mounted accelerometer and reported high sensitivity on laboratory
falls.  Subsequent work introduced two-stage detectors that combine an impact
threshold with a post-impact stillness check: a fall is declared only if a large
$\mathrm{SMV}$ spike is followed by a prolonged low-activity period, filtering out
brief vigorous activities that produce high-magnitude excursions without a subsequent
motionless phase~\parencite{igual2013challenges}.  Multi-axis variants replace the scalar
magnitude with component-specific thresholds or combine vertical, horizontal, and
resultant acceleration thresholds in a logical rule~\parencite{delahoz2014survey}.

\subsection{Mathematical limitations}

The threshold decision rule $\mathbf{1}\{\mathrm{SMV}(t) > T\}$ has three
fundamental limitations that are mathematical rather than incidental.

\textbf{No probabilistic output.}  The rule produces a hard binary decision with
no associated posterior probability.  It is therefore impossible to apply a
principled recall-precision trade-off: there is no continuous operating parameter
that smoothly exchanges false negatives for false positives.  Adjusting sensitivity
requires changing the hard threshold~$T$, but this adjustment is global---it shifts
the operating point identically for all individuals and all activities.

\textbf{No principled threshold selection.}  The threshold $T$ is chosen by a
combination of empirical cross-validation and clinical convention.  There is no
statistical framework for selecting $T$ that is optimal under a specified loss
function.  In particular, the $F_2$-optimal threshold---which depends on the
class-imbalance ratio in a given deployment population---cannot be identified from
first principles without a probabilistic model.

\textbf{No subject-specific adaptation.}  A single threshold $T$ is applied to all
individuals.  An active person with a vigorous gait produces baseline $\mathrm{SMV}$
values that approach the threshold during normal activity, generating chronic false
alarms.  A frail elderly person with a low-energy fall may not reach the threshold
at all.  The threshold detector has no mechanism to adjust to individual motion
dynamics because it has no subject-level probabilistic model.

\subsection{Comparison with the present work}

The present thesis replaces the scalar hard decision with a calibrated posterior
probability $f_\theta(\phi_\lambda(w)) \in [0,1]$.  The decision threshold $\tau$
is then a continuous parameter that can be set to maximize $F_2$ on a per-subject
basis.  The threshold sweep over
$\mathcal{T}=\{0.05,0.10,\ldots,0.95\}$ identifies the subject-optimal $\tau^*(s)$
retrospectively, and the gains---up to $+0.12$ units of $F_2$ on individual
subjects---demonstrate that subject-specific calibration provides value that is
structurally inaccessible to a global threshold detector.

% ─────────────────────────────────────────────────────────────
\section{Statistical Feature Methods}
\label{subsec:rw-statistical}

\subsection{Method family}

The dominant paradigm in wearable fall detection from the mid-2000s through the
mid-2010s combined handcrafted feature extraction with a standard supervised
classifier.  A window $w$ of length $L$ is mapped to a fixed-dimensional vector
$\phi_{\mathrm{hand}}(w)\in\R^d$ by computing scalar statistics of the raw signal
or its frequency-domain representation.  Typical time-domain features include the
window mean, variance, range, peak value, root mean square, and signal-magnitude
area; typical frequency-domain features include spectral energies in low-frequency
bands, peak frequency, and spectral entropy~\parencite{igual2013challenges,
delahoz2014survey}.  These features are fed to a classifier---most commonly
support vector machines, $k$-nearest neighbours, or decision trees---trained on
labeled windows from a fixed dataset.

Surveys by Igual et al.~\parencite{igual2013challenges} and Delahoz and
Labrador~\parencite{delahoz2014survey} catalogue dozens of systems built on this
template and report accuracy figures in the range $90\%$--$99\%$ on controlled
datasets.  These figures are broadly consistent across methods, suggesting that
the feature choice matters less than the evaluation protocol: studies reporting
the highest scores tend to use random train-test splits ignoring subject membership,
which inflates performance by allowing the model to memorize individual quirks
rather than generalize across persons.

\subsection{Mathematical limitations}

\textbf{Absence of a stability certificate.}  The relationship between the
input perturbation $\|w - w'\|$ and the feature perturbation
$\|\phi_{\mathrm{hand}}(w) - \phi_{\mathrm{hand}}(w')\|$ is not bounded in any
useful way for handcrafted statistics.  A small noise perturbation of a signal
window can produce an arbitrarily large change in a feature that includes a
zero-crossing count, an extremum detection, or a threshold-based spectral band
energy.  This instability is not a problem when training and test data come from
the same sensor and subject population, but it becomes significant under sensor
substitution, mounting-position variation, or cross-dataset evaluation.

\textbf{No geometric grounding.}  The feature set is selected by domain intuition
and validated empirically.  There is no underlying theory that explains why the
chosen statistics should separate falls from activities of daily living, nor any
guarantee that the same statistics remain discriminative when the population,
sensor, or activity set changes.  In the language of learning theory, the
inductive bias of the feature extractor is implicit and uncontrolled.

\textbf{Fixed feature dimension with no adaptation.}  The dimension $d$ is fixed by
the feature recipe and cannot be adjusted to the dataset geometry.  Different
datasets exhibit different fall dynamics---controlled laboratory falls versus
naturalistic community falls, elderly populations versus young adults---but the
feature extractor is invariant to these differences.  The hyperparameter of this
approach is the classifier, not the feature extractor itself.

\textbf{Subject-general evaluation conflates variability sources.}  Most papers in
this family report a single aggregate accuracy figure without distinguishing
within-subject repeatability from across-subject generalization.  As a consequence,
it is unclear whether the reported scores reflect a model that generalizes to new
individuals or merely one that memorizes the subjects present in the training fold.

\subsection{Comparison with the present work}

The feature map $\phi_\lambda$ of the present thesis differs from handcrafted
features in three respects.

First, the Cohen--Steiner stability theorem provides a formal Lipschitz bound:
small noise perturbations of the input window produce at most equally small
perturbations of the persistence diagram, and hence of $\phi_\lambda$.  This
guarantee is a mathematical theorem, not an empirical observation, and it holds
for any signal, any subject, and any sensor.

Second, the feature extractor is parameterized by $\lambda\in\Lambda$ and the
optimal $\lambda^*_K$ is identified by Bayesian search over a nontrivial search
space.  Different datasets converge to different configurations (SparseRips at 5~s
for MobiFall; Alpha at 1~s for SisFall; Rips at 1~s for FAD\_40Hz), demonstrating
that dataset-specific geometry is captured by the parameterisation.  No handcrafted
feature family has an analogous adaptation mechanism.

Third, the evaluation strictly separates subject-level generalization from
within-subject repeatability by using a leave-one-subject-out protocol with
$R=15$ repetitions, reporting both the mean and the 95\% confidence interval
for each combination of dataset, classifier, and protocol.

% ─────────────────────────────────────────────────────────────
\section{Deep Learning Approaches}
\label{subsec:rw-dl}

\subsection{Method family}

Convolutional neural networks (CNNs), recurrent networks (LSTMs and GRUs), and
more recently transformer-based architectures have been applied to wearable fall
detection with the aim of learning discriminative representations directly from
raw or minimally processed inertial data.  Wang et al.~\parencite{wang2020falldl}
demonstrated that a multi-scale CNN applied to three-axis accelerometer windows
achieves competitive sensitivity and specificity on MobiFall and SisFall without
manual feature engineering.  Batchuluun et al.~\parencite{batchuluun2022deep} provide
a systematic review of deep-learning fall detection systems published between 2018
and 2022 and identify CNN-LSTM hybrids and attention-based models as the current
state of the art on standard benchmarks, with reported $F_1$ scores above $0.95$
on several controlled datasets.

The apparent performance advantage of deep models over classical methods is partly
attributable to their ability to capture temporal structure across multiple scales
simultaneously, and partly to methodological differences: deep learning papers tend
to use larger datasets and more favorable evaluation protocols than their classical
counterparts, making direct comparison difficult.

\subsection{Engineering and scientific limitations}

\textbf{GPU dependency.}  Training a competitive deep network on inertial windows
requires GPU compute for gradient-based optimisation, which is practical in a
research setting but problematic for deployment on the battery-constrained processors
of wearable devices.  Inference may be feasible on a smartphone CPU with quantized
models, but training---and therefore retraining for personalisation---is not.

\textbf{Opacity.}  The features learned by a deep network's intermediate layers are
not interpretable in terms of the physical signal.  There is no explanation for why
a particular window is classified as a fall that a clinician or system designer can
evaluate independently of the training data.  TDA features, by contrast, have a direct
geometric interpretation: H$_0$ statistics measure the spread of the delay-embedded
cloud (cluster count and birth-death lifetime distribution), and H$_1$ statistics
measure the presence of loops (periodic or oscillatory structure).

\textbf{No formal stability guarantee.}  Adversarial-robustness research has
repeatedly demonstrated that small, human-imperceptible perturbations of the input
can reverse the output of a deep classifier.  For a safety-critical fall-detection
system, the absence of a stability certificate is a principled concern: a small
sensor hardware change, a firmware update affecting the analog front-end, or a
change in mounting position can degrade performance in ways that are unpredictable
without re-evaluation.  The Cohen--Steiner bound of TDA provides a certificate that
does not depend on the training distribution.

\textbf{The ``trains once, deploys forever'' paradigm.}  Deep models are trained on
a fixed corpus and deployed without adaptation.  Subject-specific personalisation
requires either fine-tuning---which demands labeled falls from the target
individual---or meta-learning frameworks that introduce significant additional
complexity and still require some labeled falls at deployment time.  In the present
thesis, subject-level adaptation is achieved by a threshold sweep that requires no
labeled falls, no model parameters to update, and no access to the training corpus.

\textbf{Computational cost at scale.}  The pipeline introduced in this thesis
consumed approximately 75~minutes of CPU time for the Bayesian representation search
on a 40-core server.  This cost is paid once per dataset, not per user.
Training a competitive deep model requires comparable or greater wall-clock time
and additionally requires GPU hardware.  For deployments targeting resource-limited
clinical or community settings, the CPU-only profile of the TDA pipeline is a
meaningful engineering advantage.

\subsection{Comparison with the present work}

The present thesis does not claim to match or exceed state-of-the-art deep learning
scores on fall detection benchmarks.  The purpose is different: to establish whether
a mathematically well-founded, computationally lightweight, and interpretable
pipeline can achieve strong $F_2$ scores under a rigorous subject-aware evaluation.
The answer is affirmative for two of the three datasets ($F_2 \geq 0.93$ on SisFall
and FAD\_40Hz) and satisfactory for the third ($F_2 \geq 0.88$ on MobiFall).  The
TDA pipeline achieves these results on a CPU without GPU access, with formally
certified feature stability, and with an interpretable feature space---properties
that deep models do not offer jointly.

% ─────────────────────────────────────────────────────────────
\section{Topological Data Analysis for Time-Series Classification}
\label{subsec:rw-tda}

\subsection{Persistent homology in applied settings}

The use of persistent homology as a feature extractor for time-series classification
has grown substantially since the early 2010s, following the development of
computationally efficient boundary matrix reduction
algorithms~\parencite{edelsbrunner2002topological} and their implementation in
open-source libraries such as GUDHI~\parencite{gudhi} and Ripser.

Pereira and de Mello~\parencite{pereira2015persistent} applied persistence diagrams to
time-series classification in several benchmark domains and showed that topological
features are competitive with or superior to classical statistical features when the
discriminative signal is encoded in the \emph{shape} of short segments rather than
their amplitude or frequency content.  Umeda~\parencite{umeda2017time} extended this line
to multivariate time series and demonstrated that the combination of $H_0$ and $H_1$
features provides complementary information not captured by either alone.

Gidea and Katz~\parencite{gidea2018topological} applied TDA to financial time series,
showing that topological features of sliding-window point clouds detect market
instability with a lead time not accessible to standard spectral methods.  Their
work established the delay-embedding framework---building a point cloud from lagged
copies of a scalar signal and computing its persistent homology---as a broadly
applicable TDA pipeline.  The pipeline introduced in the present thesis follows
this framework but extends it with a systematic search over the embedding and
filtration parameters rather than fixing them a priori.

\subsection{TDA for activity recognition and wearable signals}

Ott and Demir~\parencite{ott2017topological} applied TDA features to human activity
recognition (HAR) from smartphone accelerometer data.  Using a fixed sliding window
and pre-selected embedding parameters, they showed that persistence statistics
achieve accuracy comparable to handcrafted features on a standard HAR benchmark.
Xu et al.~\parencite{xu2019topological} similarly used TDA to distinguish activity classes
in a multiclass HAR setting, reporting improvements over baseline statistical
features on several benchmark datasets.

These works establish the viability of persistent homology for wearable inertial
classification, but differ from the present thesis in several ways that limit their
direct applicability to the fall detection problem.

\textbf{Objective.}  Both Ott et al.\ and Xu et al.\ optimise accuracy in a balanced
multiclass setting.  Fall detection is a binary problem with severe class imbalance,
and the appropriate objective is recall-oriented ($F_2$) rather than accuracy.
The distinction matters: a classifier that maximizes accuracy on an imbalanced dataset
can do so by predicting the majority class (non-fall) for almost all windows, achieving
high accuracy at the cost of near-zero recall.

\textbf{Evaluation protocol.}  The HAR studies cited above use random train-test
splits that do not control for subject membership, which inflates performance
estimates for any method by allowing individual-specific patterns to leak from
training to test.  The present thesis strictly enforces subject separation via
leave-one-subject-out cross-validation.

\textbf{Hyperparameter treatment.}  Prior TDA-for-wearables works fix the embedding
dimension $m$, delay $\tau$, and filtration type $c$ either by prior convention or
by a coarse manual grid.  The present thesis treats the full hyperparameter tuple
$\lambda=(w_\mathrm{sec}, m, \tau, r, c, d_\mathrm{met}, \varepsilon, q)$ as the
primary object of study and identifies the optimal tuple per dataset via 200
iterations of Bayesian search.  The finding that the optimal complex type differs
across all three datasets (SparseRips, Alpha, Rips) is a direct consequence of this
systematic search; it could not have been discovered by fixing the complex type to a
convention.

\textbf{Personalisation.}  Neither Ott et al.\ nor Xu et al.\ report subject-level
personalisation results or analyse the distribution of per-subject gains.  The
present thesis treats the distribution of the sweep gain $\Delta F_2(s)$ as a
scientific object in its own right, showing that it is right-skewed and
dataset-dependent and that its coefficient of variation exceeds $50\%$ in all cases.

% ─────────────────────────────────────────────────────────────
\section{Positioning of the Present Work}
\label{subsec:rw-positioning}

Table~\ref{tab:rw-comparison} summarises the structural differences between the
four approaches surveyed above and the present thesis across six dimensions.

\begin{table}[h!]
\centering
\caption{Structural comparison of fall detection and TDA-for-wearables approaches.
  \checkmark\ = property satisfied; $\times$ = property absent; $\sim$ = partially
  satisfied.}
\label{tab:rw-comparison}
\footnotesize
\begin{tabularx}{\textwidth}{lXXXXX}
\toprule
\textbf{Property} &
\textbf{Threshold} &
\textbf{Stat.\ features} &
\textbf{Deep learning} &
\textbf{TDA-HAR} &
\textbf{This work} \\
\midrule
Probabilistic output           & $\times$ & \checkmark & \checkmark & \checkmark & \checkmark \\
Formal stability certificate   & $\times$ & $\times$   & $\times$   & \checkmark & \checkmark \\
Dataset-specific tuning        & $\times$ & $\times$   & $\sim$     & $\times$   & \checkmark \\
Subject-aware evaluation       & $\times$ & $\sim$     & $\sim$     & $\times$   & \checkmark \\
Recall-oriented objective      & $\times$ & $\sim$     & $\sim$     & $\times$   & \checkmark \\
CPU-only deployment            & \checkmark & \checkmark & $\times$ & \checkmark & \checkmark \\
Subject-level personalisation  & $\times$ & $\times$   & $\times$   & $\times$   & \checkmark \\
\bottomrule
\end{tabularx}
\end{table}

The table makes the novelty of the present work concrete.  No prior approach combines
a formal stability certificate (Theorem), dataset-specific parameterisation (via
Bayesian search), subject-aware evaluation (LOSO with repetitions), a recall-oriented
objective ($F_2$), and subject-level personalisation (threshold sweep) within a
single unified pipeline.  The closest prior work---TDA for HAR---satisfies the
stability and CPU constraints but lacks all three evaluation properties.

The gap that this thesis directly addresses is the intersection of the last four
rows: a principled, recall-oriented, subject-aware, personalised fall detector that
carries a mathematical stability guarantee and is deployable on consumer hardware
without GPU acceleration.

